%% SCRAP: archive/legacy/phase-1/Reference/physics_experiment/PEER_REVIEW_SUBMISSION/06_ASCIIDOC/02-formal-theorems
%% SOURCE: docs/working/archive/legacy/phase-1/Reference/physics_experiment/PEER_REVIEW_SUBMISSION/06_ASCIIDOC/02-formal-theorems.adoc
%% STATUS: SUPERSEDED
%% FITS: proofs/ — the Isabelle/HOL proof targets stated here pre-date the 19-theory
%%   proof directory; assess whether the \texttt{theorem} sketches below introduce
%%   different hypotheses or naming conventions compared with the checked-in .thy files.
%% EDITORIAL: lifted — prose rewritten to press voice
%% TODO(bob): compare these theorem statements against StarForth_Correctness.thy and
%%   the loop theories to verify naming parity before promoting to proofs/.

\section{Formal Theorems --- Pre-Publication Record}

The following four theorems were stated prior to the Isabelle/HOL proof development.
They record the empirical evidence base and the intended formal targets.

\subsection{Theorem 1: Algorithmic Determinism}

\[
\forall P \in \text{ValidPrograms},\; \forall c \in \text{Configurations}:\;
\textit{execute}(P, c, t_1) = \textit{execute}(P, c, t_2)
\]

\textbf{Evidence.} Cache hit rates of 0\% coefficient of variation across 30 runs
per configuration ($p < 10^{-30}$, ruling out chance coincidence). The intended
Isabelle/HOL target:

\begin{lstlisting}[language=isabelle]
theorem physics_model_is_deterministic:
  forall P config t1 t2.
    valid_program P -->
    configuration_valid config -->
    run_program(P, config, t1) = run_program(P, config, t2)
\end{lstlisting}

\subsection{Theorem 2: Adaptive Convergence}

Let $\text{OptimalityGap}(t) = P_{\text{optimal}} - P_{\text{actual}}(t)$.
The system is convergent if $\exists\, t^*$ such that for all $t > t^*$,
$|P_{\text{actual}}(t)| < |P_{\text{actual}}(t-1)|$.

\textbf{Evidence.} Configuration C\_FULL: $-25.4\%$ improvement from runs~1--15 to
runs~16--30 ($p = 0.002$). A\_BASELINE and B\_CACHE show no statistically significant
change ($p = 0.654$ and $p = 0.917$ respectively).

\subsection{Theorem 3: Environmental Robustness}

\[
\text{Robustness} = \frac{\text{Algorithmic Variance}}{\text{Environmental Variance}}
= \frac{0\%}{70\%} = \infty
\]

\textbf{Evidence.} Runtime coefficient of variation 32--70\% (OS scheduling noise);
cache-hit coefficient of variation 0\% across all 90 runs. The two components are
decoupled.

\subsection{Theorem 4: Predictable Performance}

\[
P\!\left(|\hat{P}(c) - \mu(c)| \leq 2\sigma(c)\right) \geq 0.95
\]

\textbf{Evidence.} B\_CACHE configuration: $\mu = 7.80$~ms, $\sigma = 2.52$~ms,
$n = 30$, 95\% confidence interval $[2.76, 12.84]$~ms. Central Limit Theorem
validates the normal approximation at $n = 30$.

\subsection{Meta-Theorem}

A physics-driven runtime model can be formally verified through rigorous empirical
analysis, enabling safe deployment of adaptive optimisation in systems that require
formal guarantees. All four theorems above are jointly sufficient to establish
determinism, convergence, robustness, and predictability.
